\chapter{Conclusion}

This study presents a comprehensive optimization framework for enhancing electric vehicle charging infrastructure in the \glsxtrlong{kwc-cma}. Our \glsxtrlong{mip} successfully addresses current charging network deficiencies, achieving dramatic improvements in both \gls{l2} and \gls{l3} coverage while maintaining cost efficiency. The three-phase implementation strategy, requiring a net investment of \$3,995,000, provides a systematic approach to infrastructure enhancement while ensuring efficient resource allocation and minimal service disruptions.

Future research opportunities include the integration of real-time usage data, renewable energy sources, and emerging technologies such as wireless charging. These enhancements would further optimize the network's environmental impact and operational sustainability while preparing for future innovations in \glsxtrshort{ev} technology.

The methodology developed in this study provides valuable insights for urban planners and policymakers, offering an adaptable framework that could be extended to other metropolitan areas. As electric vehicle adoption accelerates across Canada, this systematic approach to charging infrastructure optimization will play a crucial role in supporting sustainable urban mobility and achieving national emissions reduction targets.